\documentclass[activite]{thedocument}
\startdatakeys
\source{}
\BO{2026}
\cycle{3}
\competencesDuSocle{RE4}
\objectifsApprentissage{C3C2at,C3C2au}
\motsCles{symétrie axiale, compas}
\enddatakeys
\begin{document}
    En utilisant {\bfseries uniquement le compas}, construire le symétrique de A par
    rapport à $(d)$. Même question pour B et C.\vskip50pt
    \begin{tikzpicture}[scale=1.2,rotate=30]
        \coordinate[label=below left:C](C)at(0,0);
        \coordinate[label=above left:B](B)at(1.5,4);
        \coordinate[label=below:A](A)at(-2,2);
        \coordinate(D1)at(3,-5);
        \coordinate[label=above right:$(d)$](D2)at(3,8);
        \draw(D1)--(D2);
        \draw node at(A){+};
        \draw node at(B){+};
        \draw node at(C){+};
    \end{tikzpicture}
\end{document}